\documentclass[11pt]{amsart}
\usepackage[margin=27mm]{geometry}
\usepackage{amsmath,amssymb,amsthm}
\usepackage{array,booktabs}
\usepackage[obeyspaces,spaces]{url}
\usepackage[colorlinks=true,linkcolor=blue,citecolor=blue,urlcolor=blue]{hyperref}
% The author's macros.
%
% some calligraphy
%
\newcommand\cA{{\mathcal A}}
\newcommand\cB{{\mathcal B}}
\newcommand\cC{{\mathcal C}}
\newcommand\cD{{\mathcal D}}
\newcommand\cE{{\mathcal E}}
\newcommand\cF{{\mathcal F}}
\newcommand\cG{{\mathcal G}}
\newcommand\cH{{\mathcal H}}
\newcommand\cI{{\mathcal I}}
\newcommand\cJ{{\mathcal J}}
\newcommand\cK{{\mathcal K}}
\newcommand\cL{{\mathcal L}}
\newcommand\cM{{\mathcal M}}
\newcommand\cN{{\mathcal N}}
\newcommand\cP{{\mathcal P}}
\newcommand\cQ{{\mathcal Q}}
\newcommand\cR{{\mathcal R}}
\newcommand\cS{{\mathcal S}}
\newcommand\cT{{\mathcal T}}
\newcommand\cU{{\mathcal U}}
\newcommand\cV{{\mathcal V}}

%
% fields and rings (and a semigroup)
%
\newcommand\cx{{\mathbb C}}% complexes
\newcommand\fld{{\mathbb F}}
\newcommand\flde{{\mathbb E}}
\newcommand\ints{{\mathbb Z}}
\newcommand\nn{{\mathbb N}}%non-negative integers
\newcommand\re{{\mathbb R}}%reals
\newcommand\rats{{\mathbb Q}}

\newcommand\Zu{{\mathbf u}}
\newcommand\Zv{{\mathbf v}}
\newcommand\Zw{{\mathbf w}}
\newcommand\Zy{{\mathbf y}}
\newcommand\Zz{{\mathbf z}}
\newcommand\Zx{{\mathbf x}}
\newcommand\Zj{{\mathbf j}}


%
% the really useful stuff
%
%\newcommand\opk[1]{\mathop{\hbox{\rm #1}}\nolimits} %% old version
\newcommand\opk[1]{\mathop{\mathrm{#1}}\nolimits}
\newcommand\comp[1]{{\mkern2mu\overline{\mkern-2mu#1}}}
\newcommand\diff{\mathbin{\mkern-1.5mu\setminus\mkern-1.5mu}}% for \setminus
\newcommand\gbin[2]{\genfrac{[}{]}{0pt}{}{#1}{#2}}
\newcommand\inv{^{-1}}
\newcommand\res{\mathbin{\mkern-2.0mu\restriction\mkern-2.0mu}}% ch2
\newcommand\sbs{\subseteq}
\newcommand\sps{\supseteq}
\newcommand\seq[3]{#1_{#2},\ldots,#1_{#3}}
\newcommand\supp{\opk{supp}}
\newcommand\im{\opk{im}}
\newcommand\row{\opk{row}}
\newcommand\pmat[1]{\begin{pmatrix} #1 \end{pmatrix}}
\newcommand\cprod{\mathbin{\square}}
%
% matrix theory
%
\newcommand\ip[2]{\langle#1,#2\rangle}
\newcommand\one{{\bf1}}
\newcommand\rk{\opk{rk}}
\newcommand\tr{\opk{tr}}
\newcommand\col{\opk{col}}
\newcommand\mat[3]{\mathrm{Mat}_{#1\times #2}(#3)}

%
% some group theory
%
\newcommand\aut[1]{{\rm Aut}(#1)}
\newcommand\fx[1]{{\rm fix}(#1)}% ch2
\newcommand\grp[1]{\langle #1\rangle}
\newcommand\nrml{\vartriangleleft}
\newcommand\nrmleq{\trianglelefteq}
\newcommand\alt[1]{\opk{Alt}(#1)}
\DeclareMathOperator{\Sym}{Sym}
% End of the author's macros.
\newtheorem{theorem}{Theorem}[section]
\newtheorem{lemma}[theorem]{Lemma}
\newtheorem{proposition}[theorem]{Proposition}
\newtheorem{corollary}[theorem]{Corollary}
\theoremstyle{definition}
\newtheorem{definition}[theorem]{Definition}
\newtheorem{remark}[theorem]{Remark}
\DeclareRobustCommand{\ml}[1]{\ifmmode\text{\normalfont\small\ttfamily\detokenize{#1}}\else{\normalfont\small\nolinkurl{#1}}\fi}
\newcommand{\nb}{N}
\emergencystretch=2em
\title[Regular graphs of small valency on 8 vertices]{The regular graphs of valency at most 3 on 8 vertices}
\date{28 September 2026}

\begin{document}
\raggedbottom
\begin{abstract}
Up to isomorphism, the regular graphs of valency at most 3 on 8 vertices are the empty graph,
$4K_2$, $C_8$, $C_3+C_5$, $2C_4$, $2K_4$, and five connected cubic graphs: the cube, the Wagner
graph, $K_{3,3}$ with one vertex replaced by a triangle, $K_4$ with two vertices replaced by
triangles, and two diamonds joined by a matching. This is prerequisite SE4 of
\ml{cgss/prerequisites.md}, used in the case $\mathrm E_8$ of Remark 4.14 of [CGSS]. We give a
proof by cases on a single vertex, explicit isomorphisms, invariants that separate the graphs,
a check with SageMath, and the shape of the Lean statements.
\end{abstract}
\maketitle

\section{Statement}\label{sec:statement}

All graphs are simple and finite. For a vertex $v$ of a graph $X$ we will write $\nb(v)$ for its
neighbourhood, and $u\sim w$ when $u$ and $w$ are adjacent. We will write $X+Y$ for the disjoint
union and $mX$ for $m$ disjoint copies of $X$.

We will describe each connected cubic graph on 8 vertices as the 8-cycle
$0\sim1\sim2\sim\cdots\sim7\sim0$ on the vertex set $\{0,\dots,7\}$ together with four
\emph{chords}, a perfect matching of non-consecutive pairs. We will write $C_8[P]$ for the
8-cycle with the set $P$ of chords. We will say that a cubic graph $Y$ arises from a cubic graph
$Z$ by \emph{truncating} a vertex $x$ with neighbours $y_1,y_2,y_3$ if we delete $x$ and add a
triangle $z_1z_2z_3$ with $z_i\sim y_i$.

\begin{definition}\label{def:graphs}
We will use the following cubic graphs on $\{0,\dots,7\}$.
\begin{enumerate}
\item The \emph{cube} $Q_3=C_8[\{03,16,25,47\}]$.
\item The \emph{Wagner graph} (the M\"obius ladder on 8 vertices) $M_8=C_8[\{04,15,26,37\}]$.
\item $Y_1=C_8[\{02,15,36,47\}]$, which is $K_{3,3}$ with one vertex truncated.
\item $Y_2=C_8[\{02,15,37,46\}]$, which is $K_4$ with two vertices truncated.
\item $Y_4=C_8[\{02,13,46,57\}]$, which is two copies of the diamond $K_4-e$ with the two vertices
of valency 2 of the one joined by a matching to those of the other.
\end{enumerate}
The subscript of $Y_t$ is its number of triangles. We will also use $2K_4$, with the parts
$\{0,1,2,3\}$ and $\{4,5,6,7\}$.
\end{definition}

The descriptions by truncation and by diamonds are not needed below; they are checked in
Section~\ref{sec:sage}, and we will prove everything from the chord descriptions.

\begin{theorem}\label{thm:main}
Let $X$ be a $k$-regular graph on 8 vertices, with $k\le3$. Then $X$ is isomorphic to exactly one
of the following graphs:
\begin{enumerate}
\item for $k=0$, the empty graph $\overline{K_8}$;
\item for $k=1$, $4K_2$;
\item for $k=2$, $C_8$, $C_3+C_5$ or $2C_4$; and
\item for $k=3$, $2K_4$, $Q_3$, $M_8$, $Y_1$, $Y_2$ or $Y_4$.
\end{enumerate}
\end{theorem}

Since the complement of a $k$-regular graph on 8 vertices is $(7-k)$-regular, we obtain that up to
isomorphism and complementation there are exactly 11 regular graphs on 8 vertices; these are the
11 graphs $Y$ of step $\langle1\rangle10$ of Remark 4.14, where $T(8)$ is switched with respect to
the edge set of $Y$ (computation C22).

Each graph of the list is clearly $k$-regular on 8 vertices, so Theorem~\ref{thm:main} has two
halves: every $k$-regular graph is isomorphic to one on the list (Sections~\ref{sec:small}
and~\ref{sec:cubic}), and no two graphs on the list are isomorphic (Section~\ref{sec:inv}).

\section{Valency at most 2}\label{sec:small}

For $k=0$ there is nothing to prove. For $k=1$ every vertex has exactly one neighbour, so the edges
of $X$ form a perfect matching $\{u_1w_1,\dots,u_4w_4\}$, and $u_i\mapsto 2i-2$, $w_i\mapsto
2i-1$ is an isomorphism onto $4K_2=\{01,23,45,67\}$.

\begin{lemma}\label{lem:cycle}
A connected 2-regular graph $X$ on $m$ vertices is isomorphic to $C_m$, and $m\ge3$.
\end{lemma}

\begin{proof}
Choose a vertex $v_0$ and a neighbour $v_1$ of $v_0$, and for $i\ge1$ let $v_{i+1}$ be the
neighbour of $v_i$ other than $v_{i-1}$. Let $r$ be the least index such that
$v_r\in\{v_0,\dots,v_{r-1}\}$, say $v_r=v_j$. If $j\ge1$, then $v_j$ is adjacent to $v_{j-1}$,
$v_{j+1}$ and $v_{r-1}$; these are distinct, since $r-1\ne j\pm1$ (for $r-1=j+1$ we would have
$v_r=v_{j}=v_{r-2}$, which $v_r$ is chosen to avoid), and this contradicts 2-regularity. Thus
$j=0$, and $v_0,\dots,v_{r-1}$ are distinct and form a cycle, with $r\ge3$ since $X$ is simple.
Every vertex on this cycle has both of its neighbours on it, so the cycle is a connected component
of $X$, and since $X$ is connected, it is all of $X$. Thus $r=m$, and $v_i\mapsto i$ is an
isomorphism onto $C_m$.
\end{proof}

Let $X$ be 2-regular on 8 vertices. Each connected component of $X$ is 2-regular, and so is a
cycle of length at least 3 by Lemma~\ref{lem:cycle}. The lengths form a partition of 8 into parts
of size at least 3; these are $8$, $5+3$ and $4+4$, since three parts would need at least 9
vertices. Thus $X$ is isomorphic to $C_8$, $C_3+C_5$ or $2C_4$.

\section{Valency 3}\label{sec:cubic}

Let $X$ be a cubic graph on 8 vertices. For a vertex $v$, we will write $R(v)$ for the set of the
four vertices other than $v$ that are not adjacent to $v$, and $e(v)$ for the number of edges with
both ends in $\nb(v)$; equivalently, $e(v)$ is the number of triangles through $v$. For disjoint
sets $S$ and $T$ of vertices, $e(S)$ is the number of edges inside $S$ and $e(S,T)$ the number of
edges between $S$ and $T$.

\begin{lemma}\label{lem:count}
Let $X$ be a cubic graph on 8 vertices and $v$ a vertex of $X$. Then
\[
e(\nb(v),R(v))=6-2e(v),\qquad e(R(v))=3+e(v).
\]
Each vertex of $\nb(v)$ has exactly $2-(\text{its valency in }\nb(v))$ neighbours in $R(v)$, and
each vertex of $R(v)$ has exactly $3-(\text{its valency in }R(v))$ neighbours in $\nb(v)$.
\end{lemma}

\begin{proof}
The last sentence holds since $v$ is adjacent to every vertex of $\nb(v)$ and to no vertex of
$R(v)$. Summing the valencies over $\nb(v)$ gives $9=3+2e(v)+e(\nb(v),R(v))$, and summing them
over $R(v)$ gives $12=2e(R(v))+e(\nb(v),R(v))$; the result follows.
\end{proof}

We will also use the graphs on 4 vertices with a given number of edges, which a reader can list
by hand, and a kernel can list by deciding over the 64 graphs on $\{1,2,3,4\}$.

\begin{lemma}\label{lem:four}
A graph $Z$ on 4 vertices with $6$, $5$ or $4$ edges is $K_4$, $K_4-e$, or one of $C_4$ and the
\emph{paw} (a triangle with a pendant edge), respectively. A triangle-free graph on 4 vertices
with 3 edges is $K_{1,3}$ or the path $P_4$.
\end{lemma}

\begin{proof}
With 4 edges, the complement has 2 edges, which are disjoint (and $Z=C_4$) or share a vertex
(and $Z$ is the paw). With 3 edges and no triangle: if a vertex has valency 3, then $Z=K_{1,3}$.
Otherwise every vertex has valency at most 2, so $Z$ is a disjoint union of paths and cycles; a
cycle would be a triangle or would have 4 edges, so $Z$ is a forest with 4 vertices and 3
edges, and so a tree of maximum valency 2, which is $P_4$.
\end{proof}

Let $m=\max_v e(v)\in\{0,1,2,3\}$, and choose a vertex $v$ with $e(v)=m$. Write $R=R(v)$. By
Lemma~\ref{lem:count}, $e(R)=3+m$. In each case below we will label $\nb(v)=\{x_1,x_2,x_3\}$ and
$R=\{w_1,w_2,w_3,w_4\}$ and determine all edges; the conclusion is a Hamiltonian cycle
$y_0y_1\cdots y_7$ of $X$ for which $y_i\mapsto i$ is an isomorphism onto the model, which the
reader can check by comparing the four remaining edges with the chords.

\subsection*{\texorpdfstring{Case $m=3$}{Case m=3}}
Then $\nb(v)$ is a triangle, so $\{v\}\cup\nb(v)$ spans $K_4$ and $e(\nb(v),R)=0$. Thus
$e(R)=6$ and $R$ spans $K_4$, and $X\cong2K_4$ (send $v,x_1,x_2,x_3$ to $0,1,2,3$ and
$w_1,\dots,w_4$ to $4,\dots,7$).

\subsection*{\texorpdfstring{Case $m=2$}{Case m=2}}
The two edges in $\nb(v)$ share a vertex; label them $x_1x_2$ and $x_2x_3$, so $x_1\not\sim x_3$.
Since $e(R)=5$, $R$ spans $K_4-e$ (Lemma~\ref{lem:four}); label it so that $w_1\not\sim w_2$. By
Lemma~\ref{lem:count}, $x_2$ has no neighbour in $R$, $x_1$ and $x_3$ have one each, $w_3$ and
$w_4$ have no neighbour in $\nb(v)$, and $w_1$ and $w_2$ have one each. Thus the two edges between
$\nb(v)$ and $R$ form a matching of $\{x_1,x_3\}$ with $\{w_1,w_2\}$, and after exchanging $w_1$
and $w_2$ we may assume that $x_1\sim w_1$ and $x_3\sim w_2$. All edges are now determined, and
$x_3\,v\,x_2\,x_1\,w_1\,w_3\,w_4\,w_2$ is a Hamiltonian cycle whose remaining edges $vx_1$,
$x_2x_3$, $w_1w_4$, $w_3w_2$ are the chords $13$, $20$, $46$, $57$. Thus $X\cong Y_4$.

\subsection*{\texorpdfstring{Case $m=1$}{Case m=1}}
Here every vertex $u$ satisfies $e(u)\le1$. Label $\nb(v)$ so that $x_1\sim x_2$ is its edge. By
Lemma~\ref{lem:count}, $x_1$ and $x_2$ have one neighbour in $R$ each, and $x_3$ has two. Since
$e(R)=4$, $R$ spans $C_4$ or the paw (Lemma~\ref{lem:four}).

\emph{The paw.} Label $R$ so that $w_1w_2w_3$ is the triangle and $w_3\sim w_4$. By
Lemma~\ref{lem:count}, $w_3$ has no neighbour in $\nb(v)$, $w_1$ and $w_2$ have one each, and
$w_4$ has two. If $x_3\not\sim w_4$, then $w_4\sim x_1$ and $w_4\sim x_2$, and $\nb(x_1)=
\{v,x_2,w_4\}$ contains the edges $vx_2$ and $x_2w_4$, so $e(x_1)=2$, a contradiction. Thus
$x_3\sim w_4$, and after exchanging $x_1$ and $x_2$ we may assume that the other neighbour of $w_4$
in $\nb(v)$ is $x_1$. The vertex $x_2$ and the second $R$-neighbour of $x_3$ now lie in
$\{w_1,w_2\}$, each of which has exactly one neighbour in $\nb(v)$; so after exchanging $w_1$ and
$w_2$ we may assume that $x_2\sim w_2$ and $x_3\sim w_1$. Then $v\,x_2\,x_1\,w_4\,w_3\,w_2\,w_1\,x_3$
is a Hamiltonian cycle whose remaining edges $vx_1$, $x_2w_2$, $w_4x_3$, $w_3w_1$ are the chords
$02$, $15$, $37$, $46$. Thus $X\cong Y_2$.

\emph{The 4-cycle.} Label $R$ as the cycle $w_1w_2w_3w_4$. By Lemma~\ref{lem:count}, every $w_i$
has exactly one neighbour in $\nb(v)$. The two neighbours of $x_3$ in $R$ are adjacent or opposite
on this cycle.

If they are adjacent, then after a rotation of the cycle they are $w_1$ and $w_2$. Thus $x_1$ and
$x_2$ are matched with $w_3$ and $w_4$, and after the reflection of the cycle that exchanges $w_1$
with $w_2$ and $w_3$ with $w_4$ we may assume that $x_1\sim w_4$ and $x_2\sim w_3$. Then
$x_2\,v\,x_1\,w_4\,w_1\,x_3\,w_2\,w_3$ is a Hamiltonian cycle whose remaining edges $x_2x_1$,
$vx_3$, $w_4w_3$, $w_1w_2$ are the chords $02$, $15$, $37$, $46$. Thus $X\cong Y_2$.

If they are opposite, then after a rotation they are $w_1$ and $w_3$. Thus $x_1$ and $x_2$ are
matched with $w_2$ and $w_4$, and after exchanging $x_1$ and $x_2$ we may assume that $x_1\sim w_2$
and $x_2\sim w_4$. Then $x_1\,v\,x_2\,w_4\,w_1\,x_3\,w_3\,w_2$ is a Hamiltonian cycle whose
remaining edges $x_1x_2$, $vx_3$, $w_4w_3$, $w_1w_2$ are the chords $02$, $15$, $36$, $47$. Thus
$X\cong Y_1$.

\subsection*{\texorpdfstring{Case $m=0$}{Case m=0}}
Then $X$ has no triangle. By Lemma~\ref{lem:count}, every $x_i$ has two neighbours in $R$, and
$e(R)=3$, so $R$ spans $K_{1,3}$ or $P_4$ (Lemma~\ref{lem:four}).

\emph{The star.} Label $R$ so that $w_4$ is the centre. By Lemma~\ref{lem:count}, $w_4$ has no
neighbour in $\nb(v)$ and $w_1,w_2,w_3$ have two each. Thus the edges between $\{x_1,x_2,x_3\}$ and
$\{w_1,w_2,w_3\}$ form a 2-regular bipartite graph, whose bipartite complement is a perfect
matching; after permuting $w_1,w_2,w_3$ we may assume that $x_i\not\sim w_i$ and $x_i\sim w_j$ for
$i\ne j$. Then $v\,x_1\,w_2\,x_3\,w_1\,w_4\,w_3\,x_2$ is a Hamiltonian cycle whose remaining edges
$vx_3$, $x_1w_3$, $w_2w_4$, $w_1x_2$ are the chords $03$, $16$, $25$, $47$. Thus $X\cong Q_3$.

\emph{The path.} Label $R$ as the path $w_1w_2w_3w_4$. By Lemma~\ref{lem:count}, $w_1$ and $w_4$
have two neighbours in $\nb(v)$, and $w_2$ and $w_3$ have one. Since $w_1\sim w_2$ and $X$ has no
triangle, $w_1$ and $w_2$ have no common neighbour, so their neighbours in $\nb(v)$ partition
$\nb(v)$; label $\nb(v)$ so that $w_1\sim x_1,x_2$ and $w_2\sim x_3$. The second neighbour of $x_3$
in $R$ is not $w_1$, and is not $w_3$ (since $x_3w_2w_3$ would be a triangle), so $x_3\sim w_4$. By
the same argument for $w_3\sim w_4$, the neighbours of $w_3$ and $w_4$ in $\nb(v)$ partition
$\nb(v)$; after exchanging $x_1$ and $x_2$ we may assume that $w_4\sim x_1$, and then $w_3\sim x_2$.
Then $v\,x_3\,w_2\,w_1\,x_1\,w_4\,w_3\,x_2$ is a Hamiltonian cycle whose remaining edges $vx_1$,
$x_3w_4$, $w_2w_3$, $w_1x_2$ are the chords $04$, $15$, $26$, $37$. Thus $X\cong M_8$.

This completes the proof that every cubic graph on 8 vertices is isomorphic to one of $2K_4$,
$Q_3$, $M_8$, $Y_1$, $Y_2$ and $Y_4$. We note that the graph $Y_2$ arises twice, from different
vertices; the case split is on a vertex, not on the graph, and the proof does not need to know
which case a given graph falls in. (The computation of Section~\ref{sec:sage} shows that $Q_3$ and
$M_8$ have $m=0$, $Y_1$ and $Y_2$ have $m=1$, $Y_4$ has $m=2$, and $2K_4$ has $m=3$.)

\section{Invariants}\label{sec:inv}

We will separate the graphs by two numbers that a kernel computation can evaluate on
$\{0,\dots,7\}$ by \ml{decide}: the number $t(X)$ of triangles, and the number $q(X)$ of 4-cycles
(subgraphs isomorphic to $C_4$). If $c(u,w)=\lvert\nb(u)\cap\nb(w)\rvert$, then
\[
q(X)=\frac12\sum_{\{u,w\}}\binom{c(u,w)}{2},
\]
where the sum runs over the unordered pairs of distinct vertices, since a 4-cycle
$u\,x\,w\,y$ is counted once at each of its two diagonals $\{u,w\}$ and $\{x,y\}$. The valency
separates graphs of different valency. Within a valency, the pairs $(t,q)$ are distinct:

\begin{center}
\begin{tabular}{lcccccc}
\toprule
$k=2$ & $C_8$ & $C_3+C_5$ & $2C_4$ \\
$(t,q)$ & $(0,0)$ & $(1,0)$ & $(0,2)$ \\
\midrule
$k=3$ & $2K_4$ & $Q_3$ & $M_8$ & $Y_1$ & $Y_2$ & $Y_4$\\
$(t,q)$ & $(8,6)$ & $(0,6)$ & $(0,4)$ & $(1,3)$ & $(2,2)$ & $(4,2)$\\
\bottomrule
\end{tabular}
\end{center}

These are easy to see by hand from the chord descriptions. For instance, in $M_8$ the 4-cycles are
$i\,(i{+}1)\,(i{+}5)\,(i{+}4)$ for $i=0,1,2,3$, and $M_8$ has no triangle since no chord joins two
vertices at distance 2 on the 8-cycle. Further invariants, which we do not need, are recorded in
Section~\ref{sec:sage}: $Q_3$ is the only bipartite cubic graph on the list, $2K_4$ the only
disconnected one, and the automorphism groups of $2K_4$, $Q_3$, $M_8$, $Y_1$, $Y_2$, $Y_4$ have
orders $1152$, $48$, $16$, $12$, $4$ and $16$.

\section{The computation}\label{sec:sage}

The script \ml{cgss/proofs/reg8/checks/reg8_check.sage}, run with SageMath~10, does the
following.
\begin{enumerate}
\item It builds the graphs of Definition~\ref{def:graphs} and checks that $Q_3$ is SageMath's
\ml{graphs.CubeGraph(3)}, that $M_8$ is the 8-cycle with the four long diagonals, that $Y_1$ is
$K_{3,3}$ with one vertex truncated, that $Y_2$ is $K_4$ with two vertices truncated, and that
$Y_4$ is the diamond graph of the case $m=2$.
\item For $k=0,1,2,3$ it lists the $k$-regular graphs on 8 vertices with \ml{nauty_geng("8 -dk
-Dk")}; it obtains 1, 1, 3 and 6 graphs, and checks that each is isomorphic to exactly one graph
of Theorem~\ref{thm:main}, and conversely.
\item It computes $t$, $q$, connectivity, bipartiteness, girth, diameter and the order of the
automorphism group of each graph, and checks that the pairs $(t,q)$ are distinct within each
valency.
\item For each cubic graph it computes $m=\max_v e(v)$ and the shapes of $R(v)$ for the vertices
with $e(v)=m$, which agree with Section~\ref{sec:cubic}.
\item It checks the six explicit isomorphisms of Section~\ref{sec:cubic}: each case graph, built
from the edges determined there, is cubic, and the stated Hamiltonian cycle maps it onto the
model.
\end{enumerate}
All checks pass. We note that the computation is an independent check and not a part of the proof,
which is complete without it.

\section{Lean}\label{sec:lean}
\raggedright

We will follow \ml{cgss/notation.md} \S10: the namespace \ml{SimpleGraph}, lowerCamelCase for
definitions, snake\_case theorems naming their conclusion, and \ml{nonempty_iso} for existence of
an isomorphism. This section gives the statement shapes that we wrote before the formalization;
the Lean formalization, in \ml{LineGraphs/Seidel/}, follows them, and \ml{cgss/notation.md}
(\S10, the row SE4) gives the final names. Mathlib (at commit \ml{082e2d37} of 4 September 2026,
the version pinned in \ml{lake-manifest.json}) has \ml{cycleGraph}, \ml{circulantGraph}, the
disjoint sum \ml{SimpleGraph.sum}, \ml{IsRegularOfDegree}, \ml{IsCycles} and Tutte's theorem, but no cube,
Wagner or truncated graphs, no statement that a connected 2-regular graph is a cycle, and no
Petersen matching theorem. So all the models below are new, and we define them on \ml{Fin 8} to
keep adjacency decidable:
\begin{center}
\begin{tabular}{ll}
\toprule
graph & Lean \\
\midrule
$4K_2$, $C_3+C_5$, $2C_4$, $2K_4$ & \ml{SimpleGraph.fourK2}, \ml{cycleThreeFive}, \ml{twoC4}, \ml{twoK4} \\
$C_8[P]$ & \ml{SimpleGraph.cycleWithChords (P : Finset (Sym2 (Fin 8)))}, \\
& \quad the supremum of \ml{cycleGraph 8} and \ml{fromEdgeSet P} \\
$Q_3$ & \ml{SimpleGraph.cubeGraph} \\
$M_8$ & \ml{SimpleGraph.wagnerGraph} \\
$Y_1$ & \ml{SimpleGraph.truncatedK33Graph} \\
$Y_2$ & \ml{SimpleGraph.twiceTruncatedK4Graph} \\
$Y_4$ & \ml{SimpleGraph.twoDiamondsGraph} \\
$q(X)$ & \ml{SimpleGraph.numFourCycles} \\
\bottomrule
\end{tabular}
\end{center}
The docstrings record the chords and the other descriptions. We will state the theorems for
\ml{X : SimpleGraph V} with \ml{[Fintype V] [DecidableEq V] [DecidableRel X.Adj]} and
\ml{hV : Fintype.card V = 8}. They are
\begin{verbatim}
theorem SimpleGraph.IsRegularOfDegree.eq_bot_of_zero
    (h : X.IsRegularOfDegree 0) : X = bot
theorem SimpleGraph.IsRegularOfDegree.nonempty_iso_fourK2
    (hV : Fintype.card V = 8) (h : X.IsRegularOfDegree 1) :
    Nonempty (Iso X fourK2)
theorem SimpleGraph.IsRegularOfDegree.nonempty_iso_cycle
    (hc : X.Connected) (h : X.IsRegularOfDegree 2) :
    Nonempty (Iso X (cycleGraph (Fintype.card V)))
\end{verbatim}
(with \ml{bot} standing for Mathlib's bottom graph, and the last one, Lemma~\ref{lem:cycle}, for
any finite $V$), together with the following.
\begin{itemize}
\item \ml{SimpleGraph.IsRegularOfDegree.nonempty_iso_of_two_of_card_eq_eight}, with conclusion the
disjunction of \ml{Nonempty (Iso X (cycleGraph 8))}, \ml{Nonempty (Iso X cycleThreeFive)} and
\ml{Nonempty (Iso X twoC4)};
\item \ml{SimpleGraph.IsRegularOfDegree.nonempty_iso_of_three_of_card_eq_eight}, with conclusion the
disjunction over \ml{twoK4}, \ml{cubeGraph}, \ml{wagnerGraph}, \ml{truncatedK33Graph},
\ml{twiceTruncatedK4Graph} and \ml{twoDiamondsGraph};
\item the helpers \ml{SimpleGraph.IsRegularOfDegree.card_edges_neighborFinset_compl} and
\ml{card_edges_compl_neighborFinset} for Lemma~\ref{lem:count}, and the four-vertex list of
Lemma~\ref{lem:four} by \ml{decide} on \ml{SimpleGraph (Fin 4)};
\item the invariants: \ml{(cubeGraph.cliqueFinset 3).card = 0}, \ml{cubeGraph.numFourCycles = 6},
and so on for the table of Section~\ref{sec:inv}, each by \ml{decide}; and from these
\ml{SimpleGraph.cubic_eight_pairwise_not_iso} and \ml{two_regular_eight_pairwise_not_iso}, via
the invariance of both counts under isomorphism (\ml{SimpleGraph.Iso}).
\end{itemize}

\emph{Why this route.} We chose the case split on a vertex $v$ maximizing $e(v)$ because every step
is local: one counting lemma, a list of graphs on 4 vertices, and at most two exchanges of labels
per case, so that a Lean proof chooses $v$, names $x_1,x_2,x_3,w_1,\dots,w_4$ by
\ml{Finset} choices, and ends by writing down the Hamiltonian cycle as an explicit
\ml{Fin 8}-valued bijection whose adjacency the kernel checks by \ml{decide}. The alternatives are
longer. A route through a perfect matching needs Petersen's theorem, which Mathlib does not have
(it has Tutte's theorem, from which Petersen's theorem is a further argument about bridgeless
graphs), and after removing the matching it leaves $C_8$, $C_3+C_5$ or $2C_4$ together with a
choice of matching, which is more cases. The complement is 4-regular with 16 edges, and its local
structure is richer. A brute-force \ml{decide} over graphs on \ml{Fin 8} is out of reach: even after
fixing $\nb(0)=\{1,2,3\}$, there remain $2^{21}$ edge sets on the other pairs.

\begin{thebibliography}{9}
\bibitem{CGSS} P.~J. Cameron, J.-M. Goethals, J.~J. Seidel and E.~E. Shult, Line graphs, root
systems, and elliptic geometry, \emph{J. Algebra} \textbf{43} (1976), 305--327.
\end{thebibliography}

\end{document}
